\chapter{Introduction} \label{chapter:intro}

Metaheuristic algorithms have proven to be powerful tools for solving complex optimization problems, especially when exact methods become computationally intractable \parencite{duan2023applications, agung2021performance, jaber2024review}. However, their performance strongly depends on the appropriate configuration of internal parameters and mechanisms \parencite{hamadi2010autonomous, dorigo1996ant, kennedy1995particle}. Achieving a balance between exploration (searching new regions) and exploitation (refining known promising areas) often requires configuration strategies that are either problem-specific or computationally expensive \parencite{huang2019survey, caceres2017experimental}. In this context, adaptive mechanisms that adjust metaheuristic behavior during execution (based on real-time feedback) have gained increasing attention for their potential to improve solution quality while reducing computational cost \parencite{lin2020can, eftimov2020dsctool}.

Fuzzy logic provides a mathematical framework for modeling imprecise information and reasoning under uncertainty, making it particularly suitable for problems where crisp thresholds and rigid rules are insufficient \parencite{zadeh1983role, zadeh1992knowledge}. Fuzzy control systems, especially those based on the Mamdani model, have been successfully applied to dynamically adjust parameters in metaheuristic algorithms \parencite{valdez2014survey, diaz2017new}. By encoding expert knowledge or behavioral patterns through linguistic rules and membership functions, fuzzy controllers allow for nuanced, real-time adjustments that improve the balance between exploration and exploitation \parencite{valdez2010fuzzy, sadollah2018introductory}. This adaptability has shown promising results in enhancing the performance of metaheuristics across a variety of problem domains \parencite{valdez2014survey, sadollah2018introductory}.

This proposal aims to develop and validate an adaptive fuzzy control system for the dynamic configuration of metaheuristics, focusing on the selection of fuzzy schemes during the execution of the optimization process. This mechanism is designed to improve the system's capacity to adapt to changing search conditions inherent to the stochastic nature of metaheuristics, thereby enabling a flexible semantic representation of the fuzzy controller’s input variables, and allowing for semantic flexibility and behavioral plasticity as the system modifies its interpretation of the optimization landscape in response to real-time indicators \parencite{valdez2014survey, dombi2021modifiers}.

Unlike traditional approaches that rely on fixed membership functions or static fuzzy controllers, this research proposes an adaptive architecture capable of dynamically selecting among a set of predefined fuzzy schemes. Each profile defines a specific fuzzy partition of the domain of a linguistic variable, consisting of a coherent set of membership functions that may differ in type, quantity, shape, distribution over the domain, and degree of overlap \parencite{sadollah2018introductory, pedrycz2015designing}. These differences allow the system to represent alternative semantic interpretations of input variables. The selection of an appropriate fuzzy profile during runtime is guided by the observed behavior of the optimization algorithm, such as stagnation, relative improvement, or diversity loss \parencite{castillo2017dynamic, patel2022type}, thereby introducing a novel form of semantic adaptation within fuzzy control systems. This capability is positioned as a methodological innovation in the dynamic configuration of metaheuristics regulating parameter control behaviour, where the output of the fuzzy system corresponds to the adaptation of metaheuristics parameters \parencite{valdez2014survey, diaz2017new}.

The dynamic selection mechanism of fuzzy schemes is not predetermined in this proposal but rather defined as a core research objective. The project will explore, implement, and comparatively evaluate several dynamic selection strategies, including static mechanisms \parencite{cord2001genetic}, stochastic approaches, reinforcement learning techniques, and hyperheuristic methods \parencite{qiu2020event, zhu2020event, haighton2024adaptable, burke2010classification, crawford2012dynamic}. The performance of these strategies will be empirically assessed to determine which is most effective in the context of adaptive metaheuristic configuration for optimization problems.

The adaptive fuzzy control system will be integrated into leading state-of-the-art metaheuristics, such as Particle Swarm Optimization (PSO) and Genetic Algorithm (GA) \parencite{valdez2014survey}, and validated through their application to classical benchmark problems, including mathematical functions from the IEEE Congress on Evolutionary Computation (CEC), combinatorial problems such as the Set Covering Problem (SCP) and the Knapsack Problem (KP) \parencite{fatih2006particle, chatterjee2024solving}, and the Feature Selection Problem \parencite{sadollah2018introductory}. A comparative experimental evaluation will be conducted to assess the performance of the different dynamic selection strategies under controlled conditions. This validation will allow for a detailed analysis of the impact of each mechanism on the effectiveness, efficiency, and robustness of the optimized algorithms \parencite{halim2021performance, djebedjian2021global, crawford2021q}. The expected outcome is a validated adaptive fuzzy control system capable of improving the performance of metaheuristic algorithms across diverse problem types.


The remainder of this document is organized into four main chapters. Chapter~\ref{chapter:marco_teorico} presents the theoretical foundations of the proposal, including the fundamentals of optimization, metaheuristics, fuzzy logic, fuzzy control systems, and the proposed semantic adaptation mechanism via fuzzy schemes. Chapter~\ref{chapter:proposal} describes the research proposal in detail, outlining the motivation, design principles, functional architecture, and validation strategy for the adaptive fuzzy control system. Chapter~\ref{chapter:marco_metodologico} defines the methodological framework, including the research questions, hypothesis, and general objective guiding the investigation. Finally, Chapter~\ref{chapter:work_plan} presents the work plan, structured into work packages with defined tasks, expected outputs, and a Gantt chart that delineates the timeline for the implementation and evaluation of the proposed system.